(1) 记 \(k_0=\omega/c\)，则
\[
k_x=n_1k_0\sin\theta,\quad p=n_1k_0\cos\theta>0,\quad
\kappa=k_0\sqrt{n_1^2\sin^2\theta-n_2^2}>0.
\]
低折射层法向波矢为 \(i\kappa\)，其单向分量呈倏逝衰减。
(2) TE场及导数在两个界面连续，以传输矩阵计算可得
\[
\mathcal T(d)=\left[1+{(p^2+\kappa^2)^2\over4p^2\kappa^2}
\sinh^2(\kappa d)\right]^{-1}.
\]
这是两端介质相同、无吸收、相同TE波阻抗时的功率透射率；用单层转移矩阵亦可复算。
(3) 两侧传播介质相同、无损且不存在辐射到其他通道，Wronskian（或法向Poynting通量）恒定，因此 \(\mathcal R+\mathcal T=1\)。\(d\to0\) 时 \(\sinh\kappa d\to0\)，\(\mathcal T\to1\)；\(d\to\infty\) 时
\[
\mathcal T\sim{16p^2\kappa^2\over(p^2+\kappa^2)^2}e^{-2\kappa d}.
\]
(4) 气隙内左右倏逝分量叠加仍可形成非零法向平均Poynting通量，但无源结构的透射功率从不超过入射；隧穿相位时间不能直接当作信号前沿速度。吸收使R+T<1，有限角谱产生束形变化，临界角附近 \(\kappa\to0\) 必须保留完整而非厚隙渐近式。
